\section{Requirements Analysis, Specification and Validation}
\label{sec:pa-analysis}

\subsection{Requirements Analysis}
\label{sec:pa-analyst}

The analyst is responsible for turning the interviews into a \emph{requirements analysis}. It receives the plan topics, the personas, the transcripts and, in a rework round, the current analysis. It produces six kinds of items: \emph{raw requirements}, the needs that a persona stated, each linked to the persona who \emph{stated} it and \emph{evidenced by} the exact words of the interview; \emph{latent needs}, implied but not stated; \emph{analyzed requirements}, single testable statements with a type and a MoSCoW priority, \emph{derived from} raw requirements or latent needs; and \emph{assumptions}, \emph{ambiguity issues} and \emph{conflicts}. An issue is marked open only when a stakeholder can settle it, and the open issues decide whether the workflow returns to the interviews (Section~\ref{sec:pa-rework}). Existing items are never edited: a changed item is a new item that \emph{supersedes} the old one, so the history of the analysis is preserved, and every requirement remains traceable to the interview statement it comes from.

\subsection{Requirement Specification}
\label{sec:pa-specification}

The requirement engineer is responsible for the \emph{requirement specification}. It receives the analyzed requirements, the latent needs, the assumptions and the conflicts, and writes the frame of the specification (purpose, scope and definitions) and requirements that are functional, non-functional or constraints. Each requirement has measurable acceptance criteria or a fit criterion, a priority, a rationale and a verification method, and is \emph{derived from} the analyzed requirements it is built from. It uses only the items it receives, so every requirement stays traceable to the interviews, in line with the quality characteristics of ISO/IEC/IEEE 29148~\cite{iso29148}.

\subsection{Product Backlog}
\label{sec:pa-backlog}

The backlog manager is responsible for the \emph{product backlog}. It receives the specification and writes epics and user stories in the form ``As a \ldots, I want \ldots, so that \ldots'', with acceptance criteria in Given/When/Then form, a priority aligned with the source requirement and an estimate in story points. The backlog covers the specification in both directions: every story is \emph{derived from} at least one requirement, and every requirement is implemented by at least one story.

\subsection{Validation}
\label{sec:pa-validation}

The validator is responsible for an independent check of the specification and the backlog; it reports problems and never rewrites the artifacts. It receives the specification, the backlog and the issues of the analysis that are still open, and produces a validation report whose issues each belong to one category, carry a severity and a suggestion, and \emph{flag} the items concerned. A pass ends the run; otherwise the run returns to the requirement engineer, until the maximum review rounds are reached, after which the run closes with the last report instead of failing.
